% Abstract

Kenya's livestock sector contributes 42\% of agricultural GDP and 12\% of national GDP, yet quantitative planning remains constrained by the absence of reproducible, data-driven analytical models. This study presents a cohort-based quarterly forecasting framework spanning nine value chains: dairy cattle, beef cattle, camels, goats, sheep, poultry, pigs, rabbits, and apiculture, initialized from the 2019 KNBS livestock census with projections from 2020 to 2041. The model applies a quarterly transition sequence covering mortality, culling, offtake, age-based maturation, and reproduction across all cohorts. Production modules translate herd and flock dynamics into species-specific output metrics, while a feed demand module estimates dry matter requirements by species and agro-ecological zone. Projected outputs indicate a CAGR of 3.6-6\% across all value chains. Beyond forecasting, the framework exposes productivity bottlenecks and structural inefficiencies, including feed deficits, low offtake rates, and suboptimal reproductive performance, that aggregate assessments typically obscure. The model provides an analytical basis for food security planning, targeted infrastructure investment, and evidence-based livestock policy.